\documentclass[11pt]{article}
\usepackage[margin=1in]{geometry}
\usepackage{amsmath,amsthm,amssymb,booktabs,hyperref,xcolor,array,colortbl}

\definecolor{proved}{HTML}{1a7a1a}
\definecolor{killed}{HTML}{b30000}
\definecolor{lightgreen}{HTML}{e6f4e6}
\definecolor{lightred}{HTML}{fde8e8}

\newtheorem{theorem}{Theorem}[section]
\newtheorem{proposition}[theorem]{Proposition}
\newtheorem{remark}[theorem]{Remark}

\title{\textbf{Positive-GCD $Q$: Pair Bounds,\\
Positive Definiteness, and a Zero-Diagonal Kill}\\[4pt]
\large\textit{Closure note --- Theorems A/B/C proved; C$'$ false}}
\author{Jonathan Robert Simons\\
Savannah, Georgia, USA\\
\texttt{simonsmedicalinnovations@gmail.com}}
\date{Corrected preprint \quad|\quad August 2026}

\begin{document}
\maketitle

\begin{center}
\fbox{\parbox{0.92\linewidth}{\centering\small
\textbf{Operator.} $Q_{ij}=\gcd(i,j)/\sqrt{ij}$ (positive GCD --- not inverse-GCD $\widetilde Q$).
No Navier--Stokes claim.}}
\end{center}

\vspace{0.5em}
\begin{center}
\begin{tabular}{>{\raggedright\arraybackslash}p{9cm}>{\centering\arraybackslash}p{3cm}}
\toprule
\textbf{Claim} & \textbf{Status}\\
\midrule
\rowcolor{lightgreen}
A: $R(e_a-e_b)>0$ on $Q$ & \textcolor{proved}{\textbf{Proved}}\\
\rowcolor{lightgreen}
B: $R_{\widehat Q}(e_a-e_b)\ge -1$ & \textcolor{proved}{\textbf{Proved}}\\
\rowcolor{lightgreen}
C: $Q_N\succ 0$ (all $N$) & \textcolor{proved}{\textbf{Proved}}\\
\rowcolor{lightred}
C$'$: $\lambda_{\min}(\widehat Q)>-1/2$ for all $N$ & \textcolor{killed}{\textbf{False}}\\
\bottomrule
\end{tabular}
\end{center}

\begin{abstract}
On the positive-GCD Gram matrix $Q_N(i,j)=\gcd(i,j)/\sqrt{ij}$ we prove
restricted pair bounds (Theorems~A--B), positive definiteness of $Q_N$ for every
$N$ (Theorem~C, via congruence with the classical GCD matrix), and kill the
zero-diagonal full-spectrum claim $\lambda_{\min}(\widehat Q)>-1/2$.
Companion inverse-GCD Bridge$^*$ lives in \texttt{04\_q6\_inverse\_gcd.tex};
do not mix the two operators.
\end{abstract}

\section{Operator}

\[
Q_N(i,j)=\frac{\gcd(i,j)}{\sqrt{ij}},\qquad 1\le i,j\le N.
\]
In particular $Q_N(a,a)=1$. Write $\widehat Q_N$ for $Q_N$ with diagonal zeroed.
Let $G_N(i,j)=\gcd(i,j)$ and $D=\mathrm{diag}(1,2,\ldots,N)$, so
\[
Q_N = D^{-1/2}\,G_N\,D^{-1/2}.
\]

\section{Theorem A --- diagonal pair positivity}

\begin{theorem}[A]\label{thm:A}
For integers $a\neq b$,
\[
R(e_a-e_b)=1-\frac{\gcd(a,b)}{\sqrt{ab}}\;>\;0.
\]
\end{theorem}

\begin{proof}
Set $g=\gcd(a,b)/\sqrt{ab}$. Then $g<1$ for $a\neq b$, and
$R=1-g>0$.
\end{proof}

\section{Theorem B --- zero-diagonal pair bound}

\begin{proposition}[B]\label{prop:B}
For $a\neq b$,
\[
R_{\widehat Q}(e_a-e_b)
=-\frac{\gcd(a,b)}{\sqrt{ab}}
\;\ge\;-1.
\]
\end{proposition}

\begin{proof}
Immediate from $\gcd(a,b)/\sqrt{ab}\le 1$.
\end{proof}

\section{Theorem C --- positive definiteness}

\begin{theorem}[C]\label{thm:C}
For every $N\ge 1$, the matrix $Q_N$ is positive definite.
\end{theorem}

\begin{proof}
The classical GCD matrix $G_N=(\gcd(i,j))_{1\le i,j\le N}$ is positive
definite (Beslin--Ligh; equivalently $\det G_N=\prod_{k=1}^N\varphi(k)>0$
by Smith's formula, and the same for principal minors on initial segments).
The diagonal matrix $D^{-1/2}$ is invertible, so
$Q_N=D^{-1/2}G_N D^{-1/2}$ is congruent to $G_N$ and hence positive definite.
\end{proof}

\section{Kill C$'$ --- zero-diagonal full spectrum}

\begin{theorem}[C$'$, kill]\label{thm:Cp}
There exist $N$ such that $\lambda_{\min}(\widehat Q_N)<-1/2$.
\end{theorem}

\begin{proof}[Numeric certificate]
E.g.\ $\lambda_{\min}(\widehat Q_5)\approx -0.793<-1/2$
(\texttt{scripts/positive\_gcd\_floor\_verify.py}).
\end{proof}

\begin{remark}[Hygiene]
Inverse-GCD $\widetilde Q$ already has $\lambda_{\min}<-1/2$ for moderate $N$.
That kill does \emph{not} transfer to positive-GCD $Q$ (which is PD).
Zero-diagonalization destroys the PD property and produces the C$'$ counterexamples.
\end{remark}

\begin{thebibliography}{9}
\bibitem{BeslinLigh}
S.~Beslin and S.~Ligh,
\textit{Greatest common divisor matrices},
Linear Algebra Appl.\ \textbf{118} (1989), 69--76.

\bibitem{Smith1876}
H.J.S.~Smith,
\textit{On the value of a certain arithmetical determinant},
Proc.\ London Math.\ Soc.\ \textbf{7} (1876), 208--212.

\bibitem{ZenodoPaper1}
J.R.~Simons,
\textit{GCD Spectral Paper~1} (v1),
Zenodo \href{https://doi.org/10.5281/zenodo.20269738}{10.5281/zenodo.20269738}.
\end{thebibliography}

\end{document}
